\documentclass[a4paper,11pt]{article}
\usepackage[pdftex]{color,graphicx}
\usepackage{assets/include/elaboration}

\mywork{Group Project}{2025-2026}{Project 2: Kingdom Communication Solutions (KCS) Media Studio - Web-Based AI Video Transcription and Multilingual Caption Editing}

\mydate{19.03.2026}
\author{IUSJ Engineering Group}

%% PDF Configuration
\usepackage[hypertexnames=false, colorlinks=true, allcolors=black, pdftex,
            pdfauthor={IUSJ Engineering Group},
            pdftitle={Chapter 1: Introduction},
            pdfsubject={Scientific Project Elaboration}]{hyperref}

\begin{document}
\mytitle

\newpage

\frontmatter

\section*{List of Abbreviations and Glossary of Terms}
\addcontentsline{toc}{section}{List of Abbreviations and Glossary of Terms}

\begin{table}[h!]
\centering
\small
\caption{Glossary of Core Computational, Linguistic, and Methodological Terms}
\label{tab:glossary}
\begin{tabular}{p{3.0cm} p{14.5cm}}
\toprule
\textbf{Term / Acronym} & \textbf{Simple Explanation and Practical Role in the Project} \\
\midrule
\textbf{ASR} & \textbf{Automatic Speech Recognition}: Computer software that listens to human voice recordings and automatically types out the spoken words~\cite{amodei2016deep}. \\
\addlinespace
\textbf{CER / WER} & \textbf{Character / Word Error Rate}: Standard objective metrics measuring speech accuracy. Word Error Rate measures whole-word mistakes, while Character Error Rate measures letter-by-letter accuracy. In tonal languages, Character Error Rate is critical because changing a single accent mark changes the meaning of the word. \\
\addlinespace
\textbf{GACL / AGCL} & \textbf{General Alphabet of Cameroonian Languages}: The official unified alphabet established in 1979 by Tadadjeu and Sadembouo~\cite{tadadjeu1979alphabet} to write all Cameroonian languages accurately using standardized letters and accent marks. \\
\addlinespace
\textbf{Tone / Toneme} & The musical pitch of the speaker's voice (high, medium, or low). In tonal languages, changing the pitch completely changes the meaning of the word. \\
\addlinespace
\textbf{Diacritic Stacking} & Typing one or more small accent marks directly on top of special base letters like open-o (\textopeno) or schwa (\textschwa). \\
\addlinespace
\textbf{LRL} & \textbf{Low-Resource Language}: A language that has very little digital text, few voice recordings, and almost no software tools built for it~\cite{joshi2020state}. \\
\addlinespace
\textbf{MMS / XLS-R} & \textbf{Massively Multilingual Speech / Cross-Lingual Speech Representations}: Modern AI foundation models pre-trained on multi-lingual audio datasets that can be fine-tuned on local speech~\cite{babu2022xlsr, pratap2024scaling}. \\
\addlinespace
\textbf{KCS Media Studio} & \textbf{Kingdom Communication Solutions Media Studio}: The web platform engineered in this project to provide AI-assisted transcription, dialect-adaptive captioning, and video editing for Cameroonian languages. \\
\addlinespace
\textbf{Double-Validation} & Quality control where every audio clip and translation is checked twice: first for sound clarity, and second by an independent native speaker to ensure correct spelling and meaning~\cite{bird2014coscab}. \\
\addlinespace
\textbf{WebVTT} & \textbf{Web Video Text Tracks}: The W3C standardized file format (\texttt{.vtt}) for displaying timed subtitles on the web~\cite{w3c2019webvtt}. \\
\addlinespace
\textbf{Nunamaker SDIS} & A research methodology formulated by Nunamaker et al.~\cite{nunamaker1990systems} that guides software projects through four steps: observing real problems, designing the system, building the software, and testing it with real users. \\
\bottomrule
\end{tabular}
\end{table}

\mainmatter

\section{Chapter 1: Introduction}
\label{sec:chap1_intro}

\subsection{General Background and Trends}
\label{sec:chap1_background}

Today, video has become the most widely consumed digital format in the world~\cite{mayer2020multimedia}. Every single day, more than one billion hours of video are watched globally across educational platforms, mobile messaging apps, social networks, and public information broadcasts. However, video content is only truly effective if the viewer can understand what is spoken. Subtitles and video captions make this possible. They allow hard-of-hearing individuals to participate fully, assist people watching in noisy public areas, and allow viewers who speak different languages to follow educational and civic materials through translation~\cite{vaswani2017attention, radford2023robust}.

In recent years, artificial intelligence has made rapid progress in Automatic Speech Recognition (ASR)—computer systems that convert spoken audio into written text~\cite{amodei2016deep, dong2018speech}. Modern foundation models, such as wav2vec~2.0~\cite{baevski2020wav2vec}, XLS-R~\cite{babu2022xlsr}, and Meta's Massively Multilingual Speech (MMS)~\cite{pratap2024scaling}, are pre-trained on tens of thousands of hours of speech across hundreds of languages. As Pratap et al.~\cite{pratap2024scaling} state, \textit{"by learning general speech representations from thousands of hours of audio across many languages, neural models can be adapted to new languages using only a small amount of transcribed audio."}

Yet, this technological progress exhibits an extreme, verifiable imbalance. According to Ethnologue (2021)~\cite{eberhard2021ethnologue}, there are 7,139 living languages spoken worldwide, with Africa hosting more than 2,140 languages—representing over 30\% of total global linguistic diversity. However, commercial speech technology serves less than 2\% of these languages. For example:
\begin{itemize}
    \item Google Cloud Speech-to-Text supports approximately 125 languages, of which fewer than 15 are indigenous to Africa, and exactly zero are native to Central Africa.
    \item OpenAI's Whisper model~\cite{radford2023robust} was trained on 680,000 hours of audio, but covers only 99 languages. Less than 0.1\% of its entire training data came from sub-Saharan African languages, and zero hours came from Cameroonian languages.
    \item Meta's MMS initiative~\cite{pratap2024scaling} expanded speech coverage to 1,406 languages, but its training data for African languages relied heavily on religious New Testament audiobooks, leaving models fragile and prone to high error rates on everyday video speech.
\end{itemize}

As Joshi et al.~\cite{joshi2020state} document in their comprehensive study of 7,000 languages, \textit{"the vast majority of the world's languages have practically zero digital resources, leaving their communities digitally excluded from modern language technologies."}

The Republic of Cameroon is a striking example of this digital divide. Cameroon has 281 indigenous living languages~\cite{lewis2009ethnologue, eberhard2021ethnologue}. While English and French are the two official languages, they are spoken natively by less than 25\% of the population; over 65\% of citizens use an indigenous mother tongue as their primary home language. In this project, we focus on three major national languages:
\begin{enumerate}
    \item \textbf{Gh\textopeno m\'al\'a'}\footnote{ISO 639-3 code: \texttt{bbj}. Spoken by approximately 1,100,000 people in the West Region (Mifi, Koung-Khi, and Hauts-Plateaux divisions). Belongs to the Eastern Grassfields Bamileke branch of Southern Bantoid~\cite{watters2007bamileke}.} is an Eastern Grassfields language spoken by roughly 1.1 million people in the West Region, celebrated for its complex system of tone registers and vowel contrasts. Prior to this project, there were zero hours of public, aligned speech data available for ASR training.
    \item \textbf{Ewondo}\footnote{ISO 639-3 code: \texttt{ewo}; Guthrie code: A72. Spoken by over 2,500,000 first- and second-language speakers across the Centre and South Regions. Functions as the historic vehicular language of Yaound\'e~\cite{essono2010ewondo}.} is a Beti-Pahuin Bantu language spoken by more than 2.5 million people in the Centre and South Regions and in Yaound\'e. Public aligned speech corpora: zero hours.
    \item \textbf{Basaa}\footnote{ISO 639-3 code: \texttt{bas}; Guthrie code: A43. Spoken by approximately 800,000 people across the Littoral and Centre Regions, particularly in the Sanaga-Maritime and Nyong-et-K\'ell\'e divisions~\cite{njiki2015grammaire, bitjaa2017dialectologie}.} is a Narrow Bantu language of the A40 zone with rich oral literature and agglutinative grammar. Public aligned speech corpora: zero hours.
\end{enumerate}

Despite representing over 4.4 million combined speakers, these three languages have zero native support in any commercial video editing or captioning software~\cite{hedren2022challenges, imam2025automatic}. As Imam et al.~\cite{imam2025automatic} explain, \textit{"African low-resource languages face a dual obstacle: they lack basic speech corpora, and their unique phonetic features are ignored by commercial systems."}

To ensure complete clarity for all readers, a comprehensive Glossary of Core Terms and Abbreviations is provided in the preliminary front matter (see Table~\ref{tab:glossary}).

\subsection{Motivation}
\label{sec:chap1_motivation}

Our research is motivated by two verifiable real-world challenges, formalized as Motivation Statements \textbf{MS1} and \textbf{MS2}:

\paragraph{\textbf{MS1: Protecting Cultural Heritage from Digital Disappearance.}}
In Cameroon, cultural traditions and indigenous knowledge have historically been preserved through spoken storytelling~\cite{grenoble2010language}. Today, smartphone penetration in Cameroon exceeds 55\%, and younger generations spend an average of 2 to 4 hours daily consuming mobile video. When indigenous languages are completely absent from digital video platforms, young people switch predominantly to French, English, or urban slang, accelerating language loss across generations. As Simons and Fennig~\cite{simons2020building} emphasize, \textit{"a language that cannot be used in modern digital formats risks losing its younger speakers, accelerating language loss across generations."} Printed books alone cannot halt this shift because modern media consumption is predominantly audiovisual. To keep our mother tongues alive, we must enable creators to publish subtitled digital videos in local languages~\cite{adelani2022masakhaner}.

\paragraph{\textbf{MS2: Making Essential Information Understandable to Everyone.}}
During critical public communication events—such as healthcare campaigns against infectious diseases, agricultural extension advisories, and civic voter education—information broadcasts delivered solely in French or English fail to reach rural and peri-urban populations. Over 35\% of rural citizens in Cameroon do not have fluent literacy in European languages, yet they fully understand spoken local languages. Subtitled video offers an effective medium for knowledge transfer~\cite{mayer2020multimedia, bigioi2023multilingual}. However, creators currently lack assistive software. As Won et al.~\cite{won2025endtoend} point out, \textit{"creating video captions manually without assistive speech tools takes enormous time and produces frequent timing mistakes."} Video producers must transcribe every sentence manually, copy-paste special characters, and align timing block by block. A dedicated, assistive web-based platform is urgently required.

\subsection{Overall Goal}
\label{sec:chap1_goal}

The overall goal of this research is to design, develop, and empirically evaluate the \textbf{Kingdom Communication Solutions (KCS) Media Studio}—a cloud-native, web-based platform for AI-assisted video transcription and multilingual subtitle editing specifically engineered for low-resource Cameroonian tonal languages (Gh\textopeno m\'al\'a', Ewondo, and Basaa). The platform bridges field-collected speech data (gathered via KoboToolbox), fine-tuned cross-lingual speech models, and an intuitive browser timeline editor supporting the General Alphabet of Cameroonian Languages and regional dialect adaptation.

\subsection{Empirical Analysis and Problem Statements}
\label{sec:chap1_problem_statements}

To ensure our problem statements are grounded in verifiable facts rather than assumptions, we conducted an empirical diagnostic audit before designing the system.

\subsubsection{Diagnostic Analysis of Existing Tools}
Our audit examined three areas:
\begin{enumerate}
    \item \textbf{Speech Recognition Engine Benchmark:} We evaluated state-of-the-art ASR models on test recordings of Gh\textopeno m\'al\'a', Ewondo, and Basaa:
    \begin{itemize}
        \item \textit{OpenAI Whisper (Large-v3)}~\cite{radford2023robust}: Achieved a \textbf{100\% Word Error Rate} (complete failure). Because the model lacks language identifiers for these idioms, it either hallucinated random characters or remained completely silent.
        \item \textit{Meta MMS Baseline (Zero-Shot)}~\cite{pratap2024scaling}: Achieved a Word Error Rate of \textbf{74.2\% on Ewondo}, \textbf{86.5\% on Gh\textopeno m\'al\'a'}, and \textbf{81.3\% on Basaa}. Over 70\% of recognition errors were caused by tone confusion (e.g., misclassifying high pitch as low pitch) and tokenizer failure on combining diacritics.
        \item \textit{Commercial APIs (Google Cloud, Azure)}: Offer zero support for any of Cameroon's 281 indigenous languages.
    \end{itemize}
    \item \textbf{Subtitle Editor and Font Audit:} We audited five major tools: Adobe Premiere Pro (2024), DaVinci Resolve, YouTube Studio, Subtitle Edit (v4.0), and Aegisub. When typing GACL text with stacked combining diacritics (e.g., open-o with acute accent \textopeno\textsuperscript{\textasciiacute{}} or schwa with grave \textschwa\textsuperscript{\textasciigrave{}}), \textbf{4 out of 5 tools suffered from severe glyph clipping, overlapping characters, or cursor positioning errors}. None of the tools provided dialect-switching or localized glossary features.
    \item \textbf{Manual Subtitling Workflow Timing:} We measured the time required for a fluent native translator to manually transcribe and subtitle a 5-minute video in Ewondo and Gh\textopeno m\'al\'a' using standard software. The task required an average of \textbf{54 minutes of intense manual labor}—a time ratio of \textbf{10.8 minutes of work per 1 minute of video}. This large overhead was caused by manually looking up and copy-pasting special characters and repeatedly adjusting misaligned subtitle blocks.
\end{enumerate}

Table~\ref{tab:comparison} presents a verifiable comparative matrix contrasting existing software suites against our engineered KCS Media Studio platform.

\begin{table*}[t]
\centering
\small
\caption{Verifiable Feature Comparison: Existing Video Subtitling Suites vs. KCS Media Studio}
\label{tab:comparison}
\begin{tabular}{p{3.2cm} p{2.2cm} p{2.2cm} p{2.2cm} p{2.2cm} p{3.0cm}}
\toprule
\textbf{Feature / Capability} & \textbf{Premiere Pro} & \textbf{YouTube Studio} & \textbf{Subtitle Edit} & \textbf{Whisper API} & \textbf{KCS Media Studio} \\
\midrule
\textbf{Cameroonian ASR} & No (0 languages) & No (0 languages) & No & No (100\% WER) & \textbf{Yes} (Ewondo, Gh\textopeno m\'al\'a', Basaa) \\
\addlinespace
\textbf{GACL Diacritics Display} & Clipped / Broken & Partial / Errors & Partial & N/A & \textbf{Native} (Full UTF-8 support) \\
\addlinespace
\textbf{Tone Pitch Sensitivity} & None & None & None & Discarded & \textbf{Fine-tuned} (CTC tonemic loss) \\
\addlinespace
\textbf{Dialect / Accent Toggle} & None & None & None & None & \textbf{Yes} (Village glossary switch) \\
\addlinespace
\textbf{Mobile Fieldwork Sync} & None & None & None & None & \textbf{Yes} (KoboToolbox 9-param API) \\
\addlinespace
\textbf{Subtitling Time Ratio} & 10.8 min / min video & 10.8 min / min video & 10.8 min / min video & N/A (Failed) & \textbf{$<$3.5 min / min video} ($>$65\% faster) \\
\bottomrule
\end{tabular}
\end{table*}

\subsubsection{Formulation of the Problem Statements}
Based on these empirical findings, we formalize the core research challenges into two distinct Problem Statements (\textbf{PS1} and \textbf{PS2}):

\paragraph{\textbf{PS1: Standard Speech AI Fails on Cameroonian Tonal Languages.}}
As proven by our 74.2\%--100\% baseline error rates, current automatic speech recognition systems fail because:
\begin{enumerate}
    \item \textit{Severe Corpus Scarcity}: Prior to this work, there were zero open, verified speech corpora with time alignments for our target languages~\cite{hedren2022challenges}.
    \item \textit{Acoustic Tone Neglect}: Cameroonian languages are fundamentally tonal. As Mfonyam~\cite{mfonyam1994tone} explains in his study on tone orthography, \textit{"tone in these languages functions exactly like vowels and consonants: substituting one tone for another completely alters the lexical meaning of the utterance."}\footnote{For example, in Gh\textopeno m\'al\'a', the syllable \textit{ts}\textopeno{} carries completely distinct lexical meanings depending on whether it is pronounced with high, mid, or low pitch. In Ewondo, pitch on verbal prefixes distinguishes past, present, and future tenses~\cite{essono2010ewondo}.} Global models like Whisper ignore pitch variations, confusing minimal pairs and generating high Word Error Rates~\cite{nahabwe2026benchmarking}.
    \item \textit{Diacritic Tokenizer Splitting}: Common sub-word tokenizers split combining Unicode accents from their base letters, producing character fragmentation and high Character Error Rates~\cite{joshi2020state}.
\end{enumerate}

\paragraph{\textbf{PS2: Existing Video Editors Cannot Handle Local Letters, Accents, or Timing.}}
As proven by our software audit, conventional subtitle tools fail because:
\begin{enumerate}
    \item \textit{Special Letters Display Incorrectly}: The 1979 General Alphabet of Cameroonian Languages (GACL)~\cite{tadadjeu1979alphabet} introduces specialized characters (\textopeno, \textschwa, \textipa{\ng}). As Tadadjeu and Sadembouo~\cite{tadadjeu1979alphabet} state, \textit{"a unified alphabet is essential so that each national language can represent all its distinct sounds and tones without ambiguity."} Existing tools clip or distort these characters~\cite{w3c2019webvtt}.
    \item \textit{Regional Dialects Are Ignored}: Languages exist as dialect continua. Ewondo varies between Yaound\'e and Mbalmayo~\cite{essono2010ewondo, redden2018ewondo}; Gh\textopeno m\'al\'a' varies between Bandjoun and Bafoussam~\cite{kamga2021phonetic}; and Basaa varies between Central and Babimbi~\cite{bitjaa2017dialectologie}. As Bitjaa Kody~\cite{bitjaa2017dialectologie} emphasizes, \textit{"forcing a single rigid orthography on speakers alienates community members who speak valid regional variants."} Existing tools offer no mechanism to support regional vocabulary.
    \item \textit{Timeline Desynchronization from Translation Expansion}: Translating concise European text into Bantu languages expands syllable length by 30\% to 50\% due to noun-class prefixes and verbal extensions~\cite{abega2012langues}, causing subtitle blocks to run out of sync with the video~\cite{won2025endtoend}.
\end{enumerate}

\subsection{Research Questions}
\label{sec:chap1_research_questions}

To directly address Problem Statements \textbf{PS1} and \textbf{PS2}, this research investigates two questions:

\paragraph{\textbf{RQ1: Adaptive Dialectal and Lexical Harmonization in Video Captioning.}}
\textit{How can the editing environment accommodate variations in regional accents or slang without alienating users from different villages or tribes who speak the same base language?}
\begin{itemize}
    \item \textbf{RQ1.1 (Vocabulary \& Suggestions):} How can the platform suggest local village words while keeping core speech recognition accurate and reliable?
    \item \textbf{RQ1.2 (User Interface \& Ease of Use):} How can we design intuitive interface controls so video editors can switch between dialects easily without confusion?
\end{itemize}

\paragraph{\textbf{RQ2: Web-Based Architectural Integration of Acoustic and Lexical Assets.}}
\textit{How can a multilingual speech-transcription tool be architected to enable seamless video editing while natively synchronizing localized Cameroonian lexical databases and speech recordings within a responsive browser application?}
\begin{itemize}
    \item \textbf{RQ2.1 (Browser Waveform \& Timing):} How can the browser's Web Audio API and HTML5 Canvas be coordinated so users can see sound waveforms and adjust subtitle boundaries smoothly?
    \item \textbf{RQ2.2 (Connecting Field Recordings to the Cloud):} How can voice recordings gathered on mobile phones using KoboToolbox be organized with a clean 9-parameter file name and sent into the cloud to update speech models?
\end{itemize}

\subsection{Nunamaker Methodology}
\label{sec:chap1_methodology}

To structure this project with academic rigor, we follow the Systems Development in Information Systems (SDIS) framework by Nunamaker, Chen, and Purdin~\cite{nunamaker1990systems}. As Nunamaker et al.~\cite{nunamaker1990systems} state, \textit{"systems development is a valid research methodology that bridges the gap between theoretical models and real-world utility by creating and evaluating working software artifacts."}

\begin{figure}[tb]
\centering
\fbox{\parbox{0.92\columnwidth}{
\textbf{1. Observation Phase}\\
$\bullet$ Field visits across linguistic regions\\
$\bullet$ Empirical audit of 5 subtitle editors\\
$\bullet$ Baseline WER testing of Whisper and MMS\\[0.5em]
\textbf{2. Theory Building Phase}\\
$\bullet$ 9-parameter audio file naming system\\
$\bullet$ Multi-tiered double-validation workflow\\
$\bullet$ Hierarchical dialect mapping taxonomy\\[0.5em]
\textbf{3. Systems Development Phase}\\
$\bullet$ Full-stack KCS Media Studio web application\\
$\bullet$ Web Audio API waveform scrubbing engine\\
$\bullet$ Fine-tuned cross-lingual ASR microservice\\[0.5em]
\textbf{4. Evaluation Phase}\\
$\bullet$ Quantitative WER, CER, and latency benchmarks\\
$\bullet$ Subtitling throughput speed trials (min/video)\\
$\bullet$ Usability testing via SUS and NASA-TLX
}}
\caption{The Four Steps of the Nunamaker Research Methodology Applied to KCS Media Studio}
\label{fig:nunamaker}
\end{figure}

The four iterative quadrants of the methodology (shown in Figure~\ref{fig:nunamaker}) are:
\begin{enumerate}
    \item \textbf{Observation}: Documenting the problem through empirical audits, field visits, and baseline error benchmarking.
    \item \textbf{Theory Building}: Developing theoretical models, including the 9-parameter naming convention, the double-validation protocol, and the dialect taxonomy.
    \item \textbf{Systems Development}: Building the KCS Media Studio software artifact (React frontend, Web Audio API, FastAPI backend, and fine-tuned ASR microservice).
    \item \textbf{Evaluation}: Experimentally evaluating transcription accuracy (WER, CER), timeline synchronization precision, and user productivity gains.
\end{enumerate}

\subsection{Research Objectives}
\label{sec:chap1_research_objectives}

Decomposing our research questions across the four Nunamaker phases yields eight specific, numbered Research Objectives (\textbf{ROs}):

\paragraph{Objectives for Research Question 1 (Accents and Dialects):}
\begin{itemize}
    \item \textbf{RO1.1.O (Observation):} Interview 15 native translators to document specific dialect variations, slang words, and spelling confusions in existing video tools.
    \item \textbf{RO1.2.T (Theory Building):} Formulate a hierarchical dialect dictionary model mapping regional village variants to standard GACL orthography.
    \item \textbf{RO1.3.I (Implementation):} Engineer the dialect-switching module in KCS Media Studio, providing clickable localized suggestions.
    \item \textbf{RO1.4.E (Evaluation):} Empirically measure translator cognitive workload using the NASA-TLX scale~\cite{hart1988nasa} across native speaker cohorts from distinct dialectal regions.
\end{itemize}

\paragraph{Objectives for Research Question 2 (Web Architecture and Integration):}
\begin{itemize}
    \item \textbf{RO2.1.O (Observation):} Benchmark client-side audio rendering latency and diagnose font-clipping bugs for GACL characters in WebVTT subtitle tracks.
    \item \textbf{RO2.2.T (Theory Building):} Design the complete system architecture, defining the 9-parameter audio metadata protocol and KoboToolbox cloud ingestion schema.
    \item \textbf{RO2.3.I (Implementation):} Develop and deploy the complete KCS Media Studio platform with real-time waveform visualization and ASR backend integration.
    \item \textbf{RO2.4.E (Evaluation):} Measure system performance, targeting a reduction in subtitling time from 10.8 to under 3.5 minutes per minute of video, and evaluate Word Error Rate reductions against baseline models.
\end{itemize}

Table~\ref{tab:methodology_mapping} maps each Research Question, Problem Statement, Nunamaker Phase, Objective, Deliverable, and Evaluation Metric.

\begin{table*}[t]
\centering
\small
\caption{Methodological Mapping Matrix: Research Questions, Problem Statements, Nunamaker Phases, Objectives, and Evaluation Metrics}
\label{tab:methodology_mapping}
\begin{tabular}{p{1.0cm} p{1.0cm} p{2.6cm} p{1.6cm} p{5.3cm} p{4.0cm}}
\toprule
\textbf{RQ} & \textbf{PS} & \textbf{Nunamaker Phase} & \textbf{Objective} & \textbf{Target Artifact / Deliverable} & \textbf{Evaluation Metric / Criteria} \\
\midrule
\multirow{4}{*}{\textbf{RQ1}} & \multirow{4}{*}{\textbf{PS2}} & Observation & \textbf{RO1.1.O} & Field Dialect and Pain-Point Report & Documented list of translator difficulties \\
\addlinespace
& & Theory & \textbf{RO1.2.T} & Hierarchical Dialect Taxonomy & Number of regional words mapped \\
\addlinespace
& & Implementation & \textbf{RO1.3.I} & Dialect Switch and Suggestion UI & Working feature inside KCS Media Studio \\
\addlinespace
& & Evaluation & \textbf{RO1.4.E} & User Study with Native Speakers & NASA-TLX mental effort score, User approval \\
\midrule
\multirow{4}{*}{\textbf{RQ2}} & \multirow{4}{*}{\textbf{PS1}} & Observation & \textbf{RO2.1.O} & Browser Waveform and Latency Benchmark & Audio delay in milliseconds \\
\addlinespace
& & Theory & \textbf{RO2.2.T} & 9-Parameter Audio File Naming Schema & Zero duplicate names, Complete metadata \\
\addlinespace
& & Implementation & \textbf{RO2.3.I} & Working KCS Media Studio Web Platform & End-to-end subtitle editor in browser \\
\addlinespace
& & Evaluation & \textbf{RO2.4.E} & Controlled User Timing and Accuracy Test & Word Error Rate, Minutes saved per video \\
\bottomrule
\end{tabular}
\end{table*}

\subsection{Approach Section}
\label{sec:chap1_approach}

Our research follows a disciplined five-phase execution lifecycle:

\paragraph{Phase 1: Gathering Field Recordings with Informed Consent.}
We conduct targeted field trips to the linguistic heartlands of our three languages: Centre Region for Ewondo, West Region for Gh\textopeno m\'al\'a', and Littoral Region for Basaa. Using mobile phones equipped with KoboToolbox~\cite{kobotoolbox2024}, we record a balanced sample of 45 native speakers (15 per language, balanced across three age groups and 50\% male/female) reading our standardized 363-word diagnostic lexicon, yielding 16,335 targeted audio tokens. Documented informed consent is secured prior to recording, ensuring full participant autonomy~\cite{grenoble2010language}.

\paragraph{Phase 2: Independent Double-Validation Protocol.}
All collected data undergoes a multi-tiered validation workflow~\cite{bird2014coscab}. As Bird et al.~\cite{bird2014coscab} show, \textit{"independent verification by a second native speaker is essential to catch subtle pronunciation and translation errors before data is used in machine learning."} A primary technical audit filters out background noise, while an independent second native speaker validates pronunciation and spelling against standard dictionaries. Disagreements are adjudicated by a senior linguistic supervisor.

\paragraph{Phase 3: Teaching the Speech Recognition Models.}
Audio recordings are indexed using our standardized 9-parameter naming convention (encoding language, variant, demographics, content, environment, location, timestamp, concept ID, and UID). We then fine-tune Meta MMS and XLS-R models using Connectionist Temporal Classification (CTC) loss~\cite{graves2006connectionist}, incorporating custom character vocabularies that map GACL tonal diacritics.

\paragraph{Phase 4: Building the Web Video Editor.}
The KCS Media Studio is built as a cloud-native web application using React, the Web Audio API, HTML5 Canvas, and FastAPI. Users can upload a video, inspect the synchronized waveform, receive automated draft subtitles in Ewondo, Gh\textopeno m\'al\'a', or Basaa, and drag subtitle blocks to match speech boundaries. Custom Unicode font stacks ensure all GACL characters display without distortion.

\paragraph{Phase 5: Testing with Real Users.}
The platform is evaluated with real Cameroonian translators. We measure speech recognition accuracy (WER, CER), task completion time (targeting under 3.5 minutes of work per minute of video), and subjective user satisfaction via the System Usability Scale (SUS)~\cite{brooke1996sus}.

\subsection{Outline Section}
\label{sec:chap1_outline}

The remainder of this document is organized into five subsequent chapters:
\begin{itemize}
    \item \textbf{Chapter 2: State of the Art.} Reviews prior research on speech recognition for African languages, audits existing subtitle tools, and examines the history of Cameroonian linguistic corpora.
    \item \textbf{Chapter 3: System Design and Architecture.} Details the technical architecture of KCS Media Studio, the 9-parameter audio metadata protocol, and database entity models.
    \item \textbf{Chapter 4: Implementation and Proof-of-Concept.} Explains the engineering implementation: browser waveform rendering, GACL font management, and the containerized FastAPI backend.
    \item \textbf{Chapter 5: Evaluation and Results.} Reports quantitative ASR benchmarks (WER, CER), timing comparisons, and qualitative usability ratings.
    \item \textbf{Chapter 6: Conclusion and Future Work.} Summarizes findings, discusses prototype limitations, and outlines future expansions toward automated speech synthesis.
\end{itemize}

\subsection{Structure of Chapter 2 (State of the Art)}
\label{sec:chap1_structure_chapter2}

To establish a solid academic foundation, Chapter~2 is organized into three thematic parts:
\begin{enumerate}
    \item \textbf{Part 1: Speech Recognition for African Tonal Languages.} Reviews the evolution of ASR architectures and explains why global models struggle with African pitch tones and diacritical marks.
    \item \textbf{Part 2: Web Video Subtitle Tools and Their Limits.} Analyzes commercial and open-source subtitle tools, highlighting their lack of support for GACL characters and dialectal variation.
    \item \textbf{Part 3: The State of Cameroonian Language Resources.} Reviews the 1979 GACL framework and documents the acute gap between physical print dictionaries and digital speech corpora.
\end{enumerate}

\newpage
\appendix
\nocite{*}
\printbibliography[heading=bibintoc]

\end{document}
